\section*{Bab \ref{ch:b3:product} --- Ukuran Hasil Kali,
Fubini, Penggantian Variabel}

\begin{solution}{exo:b3:product:1}
Himpunan $\Delta$ tertutup di $\intcc01^2$, jadi bersifat Borel, dan $\mathcal
B(\intcc01^2)$ adalah aljabar-$\sigma$ hasil kalinya
(\cref{prop:b3:product:sections}(b)). Dengan berulang menurut satu arah:
\[
\int_{\intcc01}\Bigl(\int \mathbf 1_\Delta(x,y)
\,\dd\lambda(y)\Bigr)\dd\mu(x)
= \int \lambda(\{x\})\,\dd\mu(x) = 0 ;
\]
sedangkan menurut arah yang lain:
\[
\int_{\intcc01}\Bigl(\int\mathbf
1_\Delta(x,y)\,\dd\mu(x)\Bigr)\dd\lambda(y)
= \int \mu(\{y\})\,\dd\lambda(y) = \int 1\,\dd\lambda = 1 .
\]
Hipotesis yang gagal adalah keberhinggaan-$\sigma$ \omterm{def:b3:measure:measure}{ukuran} cacah
$\mu$ pada $\intcc01$ yang tak terbilang: sebab tak ada keluarga
terbilang berisi himpunan berukuran-$\mu$ berhingga yang menyelimutinya.
\end{solution}

\begin{solution}{exo:b3:product:2}
Pada $\intcc0A\times\intoo0{+\infty}$:
$\int_0^A\int_0^\infty\eu^{-xy}\abs{\sin x}\,\dd y\,\dd x =
\int_0^A\frac{\abs{\sin x}}x\dd x \leq A < \infty$ (lewat Tonelli untuk
nilai mutlaknya): sehingga Fubini berlaku, dan karena
$\int_0^\infty\eu^{-xy}\dd y = \frac1x$, maka
\[
\int_0^A\frac{\sin x}x\dd x
= \int_0^\infty\Bigl(\int_0^A\eu^{-xy}\sin x\,\dd
x\Bigr)\dd y
= \int_0^\infty \frac{1 - \eu^{-Ay}(\cos A + y\sin A)}{1 +
y^2}\,\dd y
\]
(dengan integral dalamnya: $\operatorname{Im}\int_0^A\eu^{(\iu -
y)x}\dd x$, yang dihitung langsung). Saat $A \to \infty$,
suku koreksinya terbatas oleh $\int_0^\infty\eu^{-Ay}\frac{1 +
y}{1 + y^2}\dd y \leq \frac32\int_0^\infty\eu^{-Ay}\dd y =
\frac3{2A} \to 0$; sedangkan suku utamanya adalah
$\int_0^\infty\frac{\dd y}{1+y^2} = \frac\pi2$. Karena itu
$\int_0^\infty\frac{\sin x}x\dd x = \frac\pi2$ --- yakni integral
Dirichlet lewat Fubini.
\end{solution}

\begin{solution}{exo:b3:product:3}
(a) Lewat kue berlapis (\cref{prop:b3:product:layercake}): $\int
f\,\dd\mu = \int_0^\infty\mu(f > t)\,\dd t$, dan $t \mapsto
\mu(f > t)$ bersifat tak naik. Pada $[n-1, n]$: $\mu(f \geq n)
\leq \mu(f > t) \leq \mu(f > n - 1) \leq \mu(f \geq n - 1)$;
lalu jumlahkan integralnya atas selang satuannya:
\[
\sum_{n\geq1}\mu(f \geq n) \leq \int f\,\dd\mu \leq
\sum_{n\geq1}\mu(f \geq n - 1) = \mu(f \geq 0) +
\sum_{n\geq1}\mu(f\geq n) \leq \mu(X) + \sum_{n\geq1}\mu(f\geq
n).
\]

(b) Terapkan (a) pada $\abs f$: maka keberhinggaan integralnya dan
deretnya saling setara (sebab tambahan $\mu(X)$ bersifat berhingga).
\end{solution}

\begin{solution}{exo:b3:product:4}
Karena $f(x,y) = \partial_y\bigl(\frac{y}{x^2+y^2}\bigr)$ untuk
$x \ne 0$:
\[
\int_0^1 f(x, y)\,\dd y = \frac{1}{x^2 + 1}
\ \Longrightarrow\
\int_0^1\!\!\int_0^1 f\,\dd y\,\dd x =
\int_0^1\frac{\dd x}{1 + x^2} = \frac\pi4 ;
\]
lalu menurut kesetangkupan gasal $f(y,x) = -f(x,y)$, urutan lainnya memberikan
$-\frac\pi4$. Untuk nilai mutlaknya: bila $0 < y < x$,
\[
\int_0^x f(x,y)\,\dd y = \Bigl[\frac{y}{x^2 +
y^2}\Bigr]_0^x = \frac1{2x},
\quad\text{dan } f \geq 0 \text{ di sana, sehingga}\quad
\int_0^1\!\!\int_0^1\abs f \geq \int_0^1\frac{\dd x}{2x} =
+\infty .
\]
Tak ada pertentangan dengan Fubini: sebab hipotesisnya $f \in L^1$
gagal, dan kedua integral berulangnya sekadar dua bilangan
yang berbeda.
\end{solution}

\begin{solution}{exo:b3:product:5}
(a) $(\mathbf 1_{\intcc01}*\mathbf 1_{\intcc01})(x) =
\lambda\bigl(\intcc01\cap\intcc{x-1}x\bigr)$: yakni $0$ bila $x
\notin \intcc02$, $x$ bila $0 \leq x \leq 1$, dan $2 - x$ bila $1
\leq x \leq 2$: yaitu tendanya. Mengonvolusikannya lagi memberikan gundukan
$\mathcal C^1$ yang kuadratik sepenggal-sepenggal pada $\intcc03$ (yakni
B-spline kuadratik): sehingga setiap konvolusinya menambah satu derajat kemulusan ---
yakni asas pemulusan di balik pemulus \cref{ch:b3:lp}.

(b) Jika $x \notin \overline{\operatorname{supp}f +
\operatorname{supp}g}$, maka ada bola di sekitar $x$ yang saling lepas
dengan himpunan jumlahnya; lalu untuk $y \in \operatorname{supp}g$, $x - y
\notin\operatorname{supp}f$, sehingga integrannya lenyap
secara identik: jadi $f * g = 0$ di dekat $x$.

(c) Untuk $x_n \to x$: $(f*g)(x_n) = \int
f(y)g(x_n - y)\,\dd y$; maka integrannya konvergen titik demi titik
(menurut \omterm{def:b3:topology:continuity}{kekontinuan} $g$) dan terdominasi oleh $\norm
g_\infty\,\abs f \in L^1$: sehingga kekonvergenan terdominasinya memberikan $(f*g)(x_n) \to (f*g)(x)$.
\end{solution}

\begin{solution}{exo:b3:product:6}
(a) Lewat Fubini dan induksi, dengan mengiris sepanjang koordinat
terakhirnya:
\[
\lambda_d(\Delta_d) = \int_0^1
\lambda_{d-1}\bigl((1 - t)\,\Delta_{d-1}\bigr)\,\dd t
= \lambda_{d-1}(\Delta_{d-1})\int_0^1(1 - t)^{d-1}\dd t
= \frac{\lambda_{d-1}(\Delta_{d-1})}{d},
\]
dengan memakai aturan pendilatan $\lambda_{d-1}(\rho A) =
\rho^{d-1}\lambda_{d-1}(A)$
(\cref{thm:b3:product:linearchange}); lalu dengan
$\lambda_1(\Delta_1) = 1$: volumenya $\frac1{d!}$.

(b) $v_2 = \pi/\Gamma(2) = \pi$; dan $v_3 =
\pi^{3/2}/\Gamma(\frac52) = \pi^{3/2}/(\frac32\cdot\frac12
\sqrt\pi) = \frac{4\pi}3$. Sedangkan elipsoidnya adalah $T(B(0,1))$ dengan
$T = \operatorname{diag}(a_1, \dots, a_d)$: sehingga
\cref{thm:b3:product:linearchange} memberikan volume
$v_d\prod a_i$.
\end{solution}

\begin{solution}{exo:b3:product:7}
Di $\R^2$, lewat \omterm{ex:b3:product:polar}{koordinat kutub} (\cref{ex:b3:product:polar}):
\[
\int_{B(0,1)}\frac{\dd x\dd y}{(x^2+y^2)^s}
= 2\pi\int_0^1 r^{1 - 2s}\,\dd r,
\qquad
\int_{\R^2\setminus B(0,1)} = 2\pi\int_1^\infty r^{1-2s}\dd r:
\]
yang berhingga jika dan hanya jika $1 - 2s > -1$ (yakni $s < 1$), berturut-turut $1 - 2s < -1$ (yakni $s >
1$). Di $\R^d$, hindarilah koordinat bola dengan memakai kue
berlapisnya: $\lambda_d(\{\norm x^{-2s} > t\}\cap B(0,1)) =
\lambda_d(B(0, \min(1, t^{-1/2s}))) = v_d\min(1, t^{-d/2s})$,
dan $\int_0^\infty v_d\min(1, t^{-d/(2s)})\dd t < \infty$ jika dan hanya jika
$\frac d{2s} > 1$, yakni $s < \frac d2$; sedangkan integral luarnya
konvergen jika dan hanya jika $s > \frac d2$ (lewat perhitungan yang sama pada daerah
komplemennya).
\end{solution}

\begin{solution}{exo:b3:product:8}
Lewat Tonelli (sebab integrannya positif) dan penggantian variabel
$(x, y) = \Phi(u, v) = (uv,\ u(1-v))$, yakni sebuah difeomorfisma
$\mathcal C^1$ dari $\intoo0\infty\times\intoo01$ ke kuadran
terbuka dengan
\[
\det D\Phi = \det\begin{pmatrix} v & u\\ 1 - v & -u
\end{pmatrix} = -uv - u(1 - v) = -u,
\qquad \abs{\det} = u :
\]
\[
\Gamma(p)\Gamma(q) =
\iint x^{p-1}y^{q-1}\eu^{-x-y}\dd x\,\dd y
= \iint (uv)^{p-1}\bigl(u(1{-}v)\bigr)^{q-1}\eu^{-u}\,u\,
\dd u\,\dd v
= \Gamma(p + q)\,B(p, q).
\]
Lalu substitusi $t = \sin^2\theta$ pada $B(p,q)$ memberikan
$2\int_0^{\pi/2}\sin^{2p-1}\theta\cos^{2q-1}\theta\,\dd\theta =
B(p, q)$. Untuk Wallis: $W_n =
\int_0^{\pi/2}\sin^n\theta\,\dd\theta = \frac12B\bigl(\frac{n +
1}2, \frac12\bigr) = \frac{\Gamma(\frac{n+1}2)\sqrt\pi}
{2\,\Gamma(\frac n2 + 1)}$ --- misalnya $W_{2n} =
\frac\pi2\cdot\frac{(2n)!}{4^n(n!)^2}$ dengan memakai
$\Gamma(n + \frac12) = \frac{(2n)!}{4^nn!}\sqrt\pi$.
\end{solution}

\begin{solution}{exo:b3:product:9}
Untuk indikatornya: $\int\mathbf 1_B\,\dd(T_*\mu) = T_*\mu(B) =
\mu(T^{-1}B) = \int\mathbf 1_B\circ T\,\dd\mu$; lalu kelinearannya
memperluasnya ke $g$ yang sederhana, dan kekonvergenan monotonnya ke $g \geq 0$ --- yakni rumus
pemindahannya. Ia murni formal: sebab ia menyatakan ulang integral terhadap
$T_*\mu$ tetapi tak mengatakan apa pun tentang $T_*\mu$ itu \emph{apa}. Sedangkan
isi \cref{thm:b3:product:linearchange} dan
\cref{thm:b3:product:changeofvar} adalah pengenalannya
\[
\Phi_*\bigl(\lambda_d\restriction_U\bigr) =
\abs{\det D\Phi^{-1}}\,\lambda_d\restriction_V
\quad\text{(yakni ukuran berkepadatan)},
\]
yakni sebuah perhitungan peta maju \omterm{def:b3:measure:lebesgueouter}{ukuran Lebesgue} ---
dengan masukan analitisnya berupa geometri diferensial $\Phi$,
bukan formalisme teori \omterm{def:b3:measure:measure}{ukuran}.
\end{solution}

\begin{solution}{exo:b3:product:10}
Lewat Tonelli fungsi Gaussnya terfaktorkan, sehingga dengan $G_1 =
\int_\R\eu^{-s^2}\dd s = \sqrt\pi$ dan $\int_\R
s^2\eu^{-s^2}\dd s = \frac{\sqrt\pi}2$ (lewat pengintegralan parsial):
\[
\int_{\R^d}x_1^2\,\eu^{-\norm x^2}\dd x =
\frac{\sqrt\pi}2\;\pi^{(d-1)/2} = \frac{\pi^{d/2}}2,
\qquad
\int_{\R^d}\norm x^2\eu^{-\norm x^2}\dd x =
d\cdot\frac{\pi^{d/2}}2
\]
(sebab menurut kesetangkupannya, $\norm x^2 = \sum_ix_i^2$ menyumbang $d$ suku
yang sama --- yakni pemeriksaan kekonsistenannya). Untuk \omterm{def:b3:measure:measure}{ukuran} yang ternormalkan
$\pi^{-d/2}\eu^{-\norm x^2}\dd x$, momen keduanya adalah
$\frac d2$.
\end{solution}

\begin{solution}{exo:b3:product:11}
(a) Berlaku $H = \Phi^{-1}(\intoo0\infty)$ untuk $\Phi(x, y) = f(x) -
y$ yang diiriskan dengan $\{y > 0\}$: yang \omterm{def:b3:lebesgue:measurable}{terukur}, sebab $(x, y)
\mapsto f(x)$ dan $(x,y)\mapsto y$ \omterm{def:b3:lebesgue:measurable}{terukur} pada
hasil kalinya (lewat komposisi dengan proyeksinya). Sedangkan
penampang-$x$ dari $H$ adalah $\intoo0{f(x)}$, yang berukuran $f(x)$:
sehingga Tonelli mengintegralkan penampangnya,
\[
\lambda_{d+1}(H) = \int_{\R^d}\lambda_1\bigl(\intoo0{f(x)}
\bigr)\,\dd x = \int_{\R^d}f\,\dd\lambda_d .
\]

(b) Grafiknya adalah $\{(x,y) : y \geq f(x)\} \cap \{y \leq
f(x)\}$, yang \omterm{def:b3:lebesgue:measurable}{terukur}; sedangkan penampang-$x$-nya berupa himpunan tunggal, yang
berukuran $0$: sehingga Tonelli memberikan $\lambda_{d+1}(\text{grafik}) =
\int 0 = 0$.

(c) Bola $S^{d-1}$ adalah gabungan dua grafik $y =
\pm\sqrt{1 - \abs{x'}^2}$ atas bola satuan $\R^{d-1}$
(dengan memisahkan koordinat terakhirnya): yakni gabungan dua himpunan nol menurut
(b), jadi nol.
\end{solution}

\begin{solution}{exo:b3:product:12}
Pada $\intoo01^2$, $\frac1{1 - xy} = \sum_{n\geq0}(xy)^n$ dengan
suku tak negatif: sehingga Tonelli mengizinkan pengintegralan suku demi suku,
\[
\iint\frac{\dd x\,\dd y}{1 - xy}
= \sum_{n\geq0}\Bigl(\int_0^1x^n\dd x\Bigr)
\Bigl(\int_0^1y^n\dd y\Bigr)
= \sum_{n\geq0}\frac1{(n+1)^2} = \zeta(2) .
\]
Untuk kasus berganti tandanya, $\frac1{1 + xy} =
\sum_n(-1)^n(xy)^n$ bukan deret positif; tetapi
integral deret \emph{mutlaknya} bernilai $\zeta(2) <
\infty$, sehingga Fubini (dengan keterintegralan yang kini tegak) berlaku:
$\iint\frac{\dd x\dd y}{1 + xy} =
\sum_n\frac{(-1)^n}{(n+1)^2}$. Untuk identitas deretnya:
\[
\sum_{n\geq1}\frac{(-1)^{n-1}}{n^2}
= \sum_{n\geq1}\frac1{n^2} - 2\sum_{k\geq1}\frac1{(2k)^2}
= \zeta(2) - \frac{\zeta(2)}2 = \frac{\zeta(2)}2 .
\]
Dengan $\zeta(2) = \frac{\pi^2}6$ (\cref{pb:b3:spectral:1}):
kedua integralnya bernilai $\frac{\pi^2}6$ dan $\frac{\pi^2}{12}$.
Kepositifan Tonelli adalah segalanya pada perhitungan
pertamanya --- sebab tak ada pemeriksaan keterintegralan yang perlu sebelum
menukarnya; sedangkan pada yang kedua, kepositifan deret
mutlaknyalah yang \emph{menyahihkan} keterintegralannya sehingga
Fubini boleh berjalan pada deret bertandanya.
\end{solution}

\begin{solution}{pb:b3:product:1}
\textbf{1.} Dengan $x = t + \sqrt t\,u$ (sehingga $\dd x = \sqrt
t\,\dd u$; dan $x$ menjelajah $\intoo0\infty$ saat $u$ menjelajah
$\intoo{-\sqrt t}\infty$):
\[
\Gamma(t{+}1) = \int_0^\infty x^t\eu^{-x}\dd x
= t^t\eu^{-t}\sqrt t\int_{-\sqrt t}^{\infty}
\exp\Bigl(t\ln\Bigl(1 + \frac u{\sqrt t}\Bigr) - \sqrt
t\,u\Bigr)\dd u,
\]
sebab $x^t = t^t\exp\bigl(t\ln(1 + u/\sqrt t)\bigr)$ dan
$\eu^{-x} = \eu^{-t}\eu^{-\sqrt tu}$.

\textbf{2.} Untuk $u$ yang tetap dan $t \to \infty$: $t\ln(1 +
u/\sqrt t) - \sqrt tu = t\bigl(\frac u{\sqrt t} -
\frac{u^2}{2t} + o(\frac1t)\bigr) - \sqrt tu = -\frac{u^2}2 +
o(1)$: sehingga $g_t(u) \to \eu^{-u^2/2}$.

\textbf{3.} Tetapkan $\psi_1(h) = \varphi(h) + \frac{h^2}4$ pada
$\intoc{-1}1$: maka $\psi_1(0) = 0$ dan $\psi_1'(h) = \frac1{1+h} -
1 + \frac h2 = \frac{h(h-1)}{2(1+h)}$, yang $\geq 0$ pada
$\intoc{-1}0$ dan $\leq 0$ pada $\intcc01$: sehingga $\psi_1 \leq 0$,
yakni $\varphi(h) \leq -h^2/4$ di sana. Tetapkan $\psi_2(h) =
\varphi(h) + ch$ pada $\intco1\infty$ dengan $c = 1 - \ln2$: maka $\psi_2(1)
= \ln2 - 1 + c = 0$ dan $\psi_2'(h) = c - \frac h{1+h} \leq c -
\frac12 < 0$: sehingga $\varphi(h) \leq -ch$ untuk $h \geq 1$. Kini untuk $t
\geq 1$: jika $\abs u \leq \sqrt t$, maka $g_t(u) =
\eu^{t\varphi(u/\sqrt t)} \leq \eu^{-t(u/\sqrt t)^2/4} =
\eu^{-u^2/4}$; sedangkan jika $u \geq \sqrt t$, maka $t\,\varphi(u/\sqrt t) \leq
-ct\cdot\frac u{\sqrt t} = -c\sqrt t\,u \leq -cu$ (sebab $t \geq
1$), sehingga $g_t(u) \leq \eu^{-cu}$. Karena itu $g_t \leq \eu^{-u^2/4} +
\eu^{-cu}\mathbf 1_{u>0}$, yang \omterm{def:b3:lebesgue:l1}{terintegralkan} dan tak bergantung pada $t \geq
1$.

\textbf{4.} Lewat kekonvergenan terdominasi: $\int_\R g_t(u)\dd u \to
\int_\R\eu^{-u^2/2}\dd u = \sqrt{2\pi}$
(\cref{ex:b3:product:polar} ditambah penskalaan
$u\mapsto\sqrt2\,u$). Lalu dengan pertanyaan 1:
\[
\Gamma(t + 1) \sim \sqrt{2\pi t}\;\Bigl(\frac
t\eu\Bigr)^t,\qquad
n! \sim \sqrt{2\pi n}\,\Bigl(\frac n\eu\Bigr)^n .
\]

\textbf{5.} Dari \cref{exo:b3:product:8}, $W_{2n} =
\frac\pi2\,\frac{(2n)!}{4^n(n!)^2} =
\frac\pi2\,4^{-n}\binom{2n}n$. Lewat Stirling:
\[
\binom{2n}{n} = \frac{(2n)!}{(n!)^2}
\sim \frac{\sqrt{4\pi n}\,(2n/\eu)^{2n}}
{2\pi n\,(n/\eu)^{2n}} = \frac{4^n}{\sqrt{\pi n}} .
\]
Maka $W_{2n} \sim \frac12\sqrt{\pi/n}$, yang konsisten dengan
rekursi $W_n = \frac{n-1}nW_{n-2}$ (yang memaksa $W_n \sim
W_{n-2}$, dan bersama $W_nW_{n-1}\cdot n = \frac\pi2$ --- yakni
hubungan Wallis klasiknya --- memberikan $W_n \sim
\sqrt{\pi/(2n)}$; jadi kedua keasimtotikannya bersesuaian).

\textbf{6.} Berlaku $\Gamma(\frac d2 + 1) \sim \sqrt{2\pi\frac
d2}\,(\frac d{2\eu})^{d/2}$, sehingga
\[
v_d = \frac{\pi^{d/2}}{\Gamma(\frac d2 + 1)}
\sim \frac{1}{\sqrt{\pi d}}\Bigl(\frac{2\pi\eu}
d\Bigr)^{d/2} \longrightarrow 0
\]
secara adigeometri (sebab untuk $d > 2\pi\eu \approx 17$, setiap faktornya
$< 1$ dan menyusut). Secara numerik $v_1 = 2$, $v_2 \approx
3.14$, $v_3 \approx 4.19$, $v_4 \approx 4.93$, $v_5 \approx
5.26$, dan $v_6 \approx 5.17$: jadi maksimumnya di $d = 5$.

\textbf{7.} Dengan $k = \frac n2 + \frac{s\sqrt n}2$ (yang bulat,
$n$ genap, dan $s$ tetap): ambillah logaritma pada $2^{-n}\binom nk =
2^{-n}\frac{n!}{k!(n-k)!}$ lalu terapkan Stirling pada ketiga
faktorialnya. Dengan menulis $k = \frac n2(1 + \varepsilon)$ dan $n - k =
\frac n2(1 - \varepsilon)$ dengan $\varepsilon = s/\sqrt n$:
\[
\ln\Bigl(2^{-n}\binom nk\Bigr)
= -\frac n2\bigl[(1{+}\varepsilon)\ln(1{+}\varepsilon) +
(1{-}\varepsilon)\ln(1{-}\varepsilon)\bigr]
+ \frac12\ln\frac{2}{\pi n(1 - \varepsilon^2)} + o(1),
\]
sedangkan kurungnya adalah $\varepsilon^2 + O(\varepsilon^4) =
\frac{s^2}n + O(n^{-2})$: sehingga tampilannya menuju $-\frac{s^2}2 +
\frac12\ln\frac2{\pi n}$ sampai $o(1)$, yakni
\[
2^{-n}\binom nk \sim \sqrt{\frac{2}{\pi
n}}\;\eu^{-s^2/2} :
\]
yakni profil Gauss lemparan koin, yang \omterm{def:b3:lebesgue:measurable}{terukur} ---
yaitu bentuk lokal de Moivre--Laplace, yang akan diglobalkan di
\cref{ch:b3:clt}.

\textbf{8.} (i) Kekonvergenan terdominasinya mengubah limit titik demi titik dari pertanyaan 2
menjadi kekonvergenan integralnya, dengan memakai pendominasi
pertanyaan 3. (ii) Integral Gaussnya menilaikan limit
$\int\eu^{-u^2/2} = \sqrt{2\pi}$ --- sehingga konstanta Stirling
$\sqrt{2\pi}$ \emph{adalah} integral Gaussnya. (iii) Sedangkan
substitusi $x = t + \sqrt tu$ merupakan penggantian variabel
afin: yakni keinvarianan translasi dan aturan penskalaan
\omterm{def:b3:measure:lebesgueouter}{ukuran Lebesgue} (\cref{thm:b3:product:linearchange} dalam
dimensi $1$).

\textbf{9.} Uraikan kedua sukunya:
\[
d_n - d_{n+1} = \ln\frac{n!}{(n+1)!} + \Bigl(n +
\frac32\Bigr)\ln(n+1) - \Bigl(n + \frac12\Bigr)\ln n - 1
= \Bigl(n + \frac12\Bigr)\ln\frac{n+1}n - 1,
\]
sebab suku $-\ln(n+1)$ dan $(n + \frac32)\ln(n+1)$ bergabung
menjadi $(n + \frac12)\ln(n+1)$.

\textbf{10.} Untuk $t = \frac1{2n+1}$: $\frac{1+t}{1-t} =
\frac{2n+2}{2n} = \frac{n+1}n$, dan $n + \frac12 =
\frac1{2t}$; lalu deret gasal $\ln\frac{1+t}{1-t} =
2\sum_{k\geq0}\frac{t^{2k+1}}{2k+1}$ memberikan
\[
\Bigl(n + \frac12\Bigr)\ln\frac{n+1}n
= \sum_{k\geq0}\frac{t^{2k}}{2k+1}
= 1 + \frac{t^2}3 + \frac{t^4}5 + \cdots
\]
Lalu kurangkan $1$. Batas bawahnya: suku pertamanya saja,
$\frac{t^2}3 = \frac1{3(2n+1)^2}$. Batas atasnya: turunkan semua
penyebutnya menjadi $3$ lalu jumlahkan deret geometrinya:
$\frac{t^2}{3(1 - t^2)} = \frac1{3((2n+1)^2 - 1)} =
\frac1{12n(n+1)} = \frac1{12n} - \frac1{12(n+1)}$.

\textbf{11.} Dengan menjumlahkan batas atasnya dari $n$ sampai $\infty$
(memakai $d_m \to 0$): $d_n < \frac1{12n}$. Untuk batas bawahnya:
$\frac1{12m+1} - \frac1{12(m+1)+1} =
\frac{12}{(12m+1)(12m+13)}$, dan
\begin{align*}
\frac1{3(2m+1)^2} > \frac{12}{(12m+1)(12m+13)}
&\iff (12m+1)(12m+13) > 36(2m+1)^2 \\
&\iff 168m + 13 > 144m + 36,
\end{align*}
yang benar untuk $m \geq 1$. Lalu dengan menjumlahkan minoran teleskopik ini:
$d_n > \frac1{12n+1}$. Dan mengeksponensialkannya memberikan pengapitan
klasik $n!$.

\textbf{12.} (a) Galat nisbinya $= \eu^{d_n} - 1 <
\eu^{1/(12n)} - 1 < \frac{1.1}{12n}$ untuk $n$ yang besar; jadi $<
10^{-6}$ segera setelah $12n \geq 1.1\cdot10^6$, dan syarat
$n \geq 83\,334$ sudah cukup (sebab $\frac1{12n} \leq 10^{-6}$
sudah mengakibatkannya). (b) $\log_{10}(100!) =
\frac12\log_{10}(200\pi) + 200 - 100\log_{10}\eu +
d_{100}\log_{10}\eu = 1.39906 + 200 - 43.42945 + 0.00036
\approx 157.96997$, sehingga $100! \approx 10^{0.96997} \cdot
10^{157} = 9.333\cdot10^{157}$; sedangkan jendela terjaminnya
$(\eu^{1/1201}, \eu^{1/1200})$ berlebar di bawah $10^{-6}$ secara
nisbi --- yakni rumus ``asimtotik'' yang, pada $n =
100$, menjadi alat presisi.

\textbf{13.} Tulis $\sin^n = \sin^{n-2} - \sin^{n-2}\cos^2$,
lalu integralkan suku keduanya secara parsial (dengan $u = \cos\theta$,
$\dd v = \sin^{n-2}\cos\theta\,\dd\theta$, dan $v =
\frac{\sin^{n-1}}{n-1}$):
\[
\int_0^{\pi/2}\sin^{n-2}\cos^2 =
\Bigl[\cos\theta\,\frac{\sin^{n-1}\theta}{n-1}\Bigr]_0^{\pi/2}
+ \frac1{n-1}\int_0^{\pi/2}\sin^n = \frac{W_n}{n-1} .
\]
Karena itu $W_n = W_{n-2} - \frac{W_n}{n-1}$, yakni $W_n =
\frac{n-1}nW_{n-2}$. Dari $W_0 = \frac\pi2$ dan $W_1 = 1$:
\[
W_{2n} = \frac{(2n-1)!!}{(2n)!!}\cdot\frac\pi2
= \frac\pi2\binom{2n}n4^{-n},
\qquad
W_{2n+1} = \frac{(2n)!!}{(2n+1)!!}
= \frac{4^n}{(2n+1)\binom{2n}n},
\]
dengan mengubah faktorial rangkapnya lewat $(2n)!! = 2^nn!$ dan
$(2n-1)!! = \frac{(2n)!}{2^nn!}$. Akhirnya $nW_nW_{n-1} =
(n-1)W_{n-1}W_{n-2}$ menurut rekursinya: jadi konstan, sama dengan
$1\cdot W_1W_0 = \frac\pi2$.

\textbf{14.} Berlaku $W_{2n+1} \leq W_{2n} \leq W_{2n-1}$ (menurut kemonotonan
titik demi titik $\sin^n$) dan $\frac{W_{2n-1}}{W_{2n+1}} =
\frac{2n+1}{2n} \to 1$, yang menjepit $\frac{W_{2n}}{W_{2n+1}} \to
1$. Digabungkan dengan $W_{2n}W_{2n+1} = \frac{\pi}{2(2n+1)}$
(pertanyaan 13): $W_{2n}^2 \sim \frac\pi{4n}$, sehingga $W_{2n} \sim
\frac12\sqrt{\frac\pi n}$ dan $\binom{2n}n4^{-n} =
\frac2\pi W_{2n} \sim \frac1{\sqrt{\pi n}}$.

\textbf{15.} Pertanyaan 9--10 tak pernah memakai nilai
konstantanya: sebab dengan $e_n = \ln n! - (n+\frac12)\ln n + n$,
selisih $e_n - e_{n+1}$ terletak di $\bigl(0, \frac1{12n} -
\frac1{12(n+1)}\bigr)$, sehingga $(e_n)$ turun sedangkan $(e_n -
\frac1{12n})$ naik: yakni barisan berdampingan, yang konvergen ke
suatu $\ell$ bersama. Karena itu $n! \sim K n^{n+1/2}\eu^{-n}$ dengan $K =
\eu^\ell$.

\textbf{16.} Dengan menyubstitusikan Stirling berkonstanta tak dikenal ke
binomial pusatnya:
\[
\binom{2n}n4^{-n}\sqrt n \sim
\frac{K\,(2n)^{2n+1/2}\eu^{-2n}}{\bigl(K\,n^{n+1/2}
\eu^{-n}\bigr)^2}\,4^{-n}\sqrt n
= \frac{2^{2n}\sqrt{2}\,K\,n^{2n+1/2}}{K^2\,n^{2n+1}}
\,4^{-n}\sqrt n = \frac{\sqrt2}{K},
\]
lalu pertanyaan 14 memaksa $\frac{\sqrt2}K = \frac1{\sqrt\pi}$:
sehingga $K = \sqrt{2\pi}$. Jadi Bagian III--IV membuktikan kembali Stirling dari
nol; dan memasukkannya ke identitas Bagian I menilaikan
$\int_\R\eu^{-u^2/2}\dd u = \sqrt{2\pi}$ tanpa \omterm{ex:b3:product:polar}{koordinat
kutub}: sehingga Wallis dan Gauss saling menopang.

\textbf{17.} Berlaku $\frac{\binom{2n}{n+j}}{\binom{2n}n} =
\frac{(n!)^2}{(n+j)!\,(n-j)!} =
\prod_{i=1}^{j}\frac{n-i+1}{n+i}$ untuk $j \geq 0$ (dan menurut
kesetangkupannya untuk $j < 0$). Dengan mengambil logaritmanya, bila $1 \leq i
\leq j \leq K\sqrt n$:
\[
\ln\frac{n-i+1}{n+i} = \ln\Bigl(1 - \frac{i-1}n\Bigr) -
\ln\Bigl(1 + \frac in\Bigr) = -\frac{2i-1}{n} +
O\Bigl(\frac{i^2}{n^2}\Bigr),
\]
dan $\sum_{i\leq j}(2i - 1) = j^2$, sedangkan galatnya berjumlah
$O(j^3/n^2) = O(n^{-1/2})$: sehingga secara seragam,
$\exp\bigl(-\frac{j^2}n + O(n^{-1/2})\bigr)$.

\textbf{18.} Berlaku $\eu^{-n}\frac{n^n}{n!} \sim \eu^{-n}
\frac{n^n}{\sqrt{2\pi n}\,n^n\eu^{-n}} =
\frac1{\sqrt{2\pi n}}$. Variabel Poisson bermean $n$ mempunyai
simpangan baku $\sqrt n$, dan $\frac1{\sqrt{2\pi n}}$ tepat
merupakan tinggi puncak Gauss $\frac1{\sigma\sqrt{2\pi}}$:
yakni teorema limit pusat lokalnya, yang dicuplik di modusnya.

\textbf{19.} Untuk $a \in \intoo01$, lema kemiringan
\cref{pb:b3:lebesgue:1} (yakni pertanyaan 14 di sana), yang diterapkan pada
$\log\Gamma$ yang cembung di sekitar $n$, memberikan $(n-1)^a \leq
\frac{\Gamma(n+a)}{\Gamma(n)} \leq n^a$: sehingga nisbahnya terhadap $n^a$
terjepit oleh $(1 - \frac1n)^a \to 1$. Sedangkan untuk $a = m + a'$
(dengan $m \in \N$ dan $a' \in \intco01$): $\Gamma(n+a) = (n + a - 1)
\cdots(n + a')\Gamma(n + a')$, dan setiap dari $m$ faktornya bernilai
$n(1 + O(\frac1n))$: lalu kalikan taksirannya.

\textbf{20.} Rekursi $v_d = \frac{2\pi}dv_{d-2}$
(dari $\Gamma(\frac d2 + 1) = \frac d2\Gamma(\frac d2)$)
memberikan, dari $v_1 = 2$ dan $v_2 = \pi$:
\[
v_3 = \frac{4\pi}3,\quad v_4 = \frac{\pi^2}2,\quad
v_5 = \frac{8\pi^2}{15},\quad v_6 = \frac{\pi^3}6,\quad
v_7 = \frac{16\pi^3}{105}.
\]
Nisbah $\frac{2\pi}d$ melampaui $1$ tepat untuk $d \leq 6$,
sehingga setiap paritasnya naik lalu turun; dan secara numerik $v_4
\approx 4.93$, $v_5 \approx 5.26$, $v_6 \approx 5.17$: jadi
maksimum keseluruhannya di $d = 5$. Untuk fungsi pembangkitnya: $v_{2k} =
\frac{\pi^k}{k!}$, sehingga $\sum_kv_{2k}x^{2k} =
\eu^{\pi x^2}$ --- yakni semua volume bola berdimensi genap yang tergulung
dalam satu eksponensial, sekaligus taksiran peluruhan adigeometri
seketika bagi $v_d$.

\textbf{21.} Stirling pada pembilang dan penyebutnya, dengan $k =
\alpha n$:
\[
\binom n{\alpha n} \sim
\frac{\sqrt{2\pi n}\,n^n}
{\sqrt{2\pi\alpha n}\,(\alpha n)^{\alpha n}\,
\sqrt{2\pi(1-\alpha)n}\,((1-\alpha)n)^{(1-\alpha)n}}
= \frac{\eu^{nH(\alpha)}}{\sqrt{2\pi\alpha(1-\alpha)n}},
\]
sebab $n^n/(\alpha n)^{\alpha n}((1-\alpha)n)^{(1-\alpha)n} =
\alpha^{-\alpha n}(1-\alpha)^{-(1-\alpha)n} =
\eu^{nH(\alpha)}$ (dengan pangkat $n$-nya saling meniadakan: $\alpha n +
(1-\alpha)n = n$), dan $\eu^{-n}$-nya meniadakan pula. Di
$\alpha = \frac12$: $H = \ln2$ dan prafaktornya adalah
$\sqrt{2/(\pi n)}$ --- yakni pertanyaan 5 lagi. Sedangkan kecekungan sejati
$H$ (sebab turunan keduanya $-\frac1{\alpha(1-\alpha)} < 0$)
menaruh maksimumnya $\ln 2$ hanya di $\alpha = \frac12$: sehingga untuk
$\alpha \neq \frac12$, $\binom n{\alpha n}2^{-n} \approx
\eu^{-n(\ln2 - H(\alpha))}$ meluruh eksponensial --- yakni
mesin kombinatorik di balik setiap pernyataan pemusatan
tentang lemparan koin.

\textbf{22.} Dari $s_{d-1} = dv_d$ dan pertanyaan 20:
\[
s_0 = 2,\ \ s_1 = 2\pi,\ \ s_2 = 4\pi,\ \ s_3 = 2\pi^2,\ \
s_4 = \frac{8\pi^2}3,\ \ s_5 = \pi^3,\ \ s_6 =
\frac{16\pi^3}{15},
\]
yang secara numerik $2,\ 6.28,\ 12.57,\ 19.74,\ 26.32,\ 31.01,\
33.07$; dan $s_7 = \frac{\pi^4}3 \approx 32.47 < s_6$: sehingga
maksimumnya adalah bola-$6$. Lalu rekursi $s_{d+1} =
(d+2)\,v_{d+2} = (d+2)\,\frac{2\pi}{d+2}\,v_d = 2\pi v_d =
\frac{2\pi}d\,s_{d-1}$ menunjukkan kenaikan yang digerakkan $\frac{2\pi}d$
dan keruntuhan adigeometri yang sama seperti pada volumenya: sehingga lewat dimensi
ketujuh, bola menyusut lenyap lebih cepat daripada barisan geometri
mana pun.

\textbf{23.} Pertanyaan 11 mengatakan persis $\frac1{12n+1} < d_n <
\frac1{12n}$, dan
\[
\frac1{12n} - \frac1{12n+1} = \frac1{12n(12n+1)} =
O\Bigl(\frac1{n^2}\Bigr),
\]
sehingga $d_n = \frac1{12n} + O(\frac1{n^2})$ dan $\eu^{d_n} = 1 +
\frac1{12n} + O(\frac1{n^2})$; lalu mengalikannya dengan
$\sqrt{2\pi n}(n/\eu)^n$ memberikan rumus yang terkoreksi. Di $n =
10$: $\sqrt{20\pi}\,(10/\eu)^{10} = 7.92665 \times 453999.3
\approx 3\,598\,696$, yang kurang $30\,104$ (dengan galat nisbi
$8.3\cdot10^{-3}$); sedangkan mengalikannya dengan $1 + \frac1{120}$ memberikan
$3\,628\,685$, yang kurang $115$ (dengan galat nisbi
$3.2\cdot10^{-5}$). Sedangkan pengapitannya sendiri memakukan $10!$ antara
$3\,598\,696\,\eu^{1/121} \approx 3\,628\,559$ dan
$3\,598\,696\,\eu^{1/120} \approx 3\,628\,808$ --- jadi batas
atasnya meleset delapan satuan pada tujuh angka.

\textbf{24.} Substitusi Bagian I $x = t + \sqrt t\,u$, yang
diterapkan pada integral terpancungnya, memberikan
\[
\int_0^{t} x^{t}\eu^{-x}\,\dd x
= t^{t}\eu^{-t}\sqrt t\int_{-\sqrt t}^{0} g_t(u)\,\dd u ,
\]
sebab rentang $0 \leq x \leq t$ menjadi $-\sqrt t \leq u \leq
0$. Pendominasi pertanyaan 3 menyelimuti $g_t\mathbf 1_{u < 0}$
pula, sehingga kekonvergenan terdominasinya melahirkan
\[
\int_{-\sqrt t}^{0}g_t(u)\,\dd u \longrightarrow
\int_{-\infty}^{0}\eu^{-u^2/2}\dd u = \frac{\sqrt{2\pi}}2 ,
\]
sedangkan pertanyaan 4 memberikan $\Gamma(t+1) \sim
t^t\eu^{-t}\sqrt t\,\sqrt{2\pi}$. Jadi nisbahnya menuju
$\frac12$. Secara peluang: sebuah peubah acak Gamma berbentuk
besar menaruh separuh massanya secara asimtotik pada masing-masing sisi
modusnya --- yakni kesetangkupan teorema limit pusatnya, yang terbaca
dari satu substitusi.

\textbf{25.} Misalkan $\lambda = \frac{\alpha}{1-\alpha} \in
\intoc01$. Untuk $k \leq \lfloor\alpha n\rfloor \leq \alpha n$
berlaku $\lambda^{k - \alpha n} \geq 1$, sehingga
\[
\sum_{k=0}^{\lfloor\alpha n\rfloor}\binom nk
\leq \lambda^{-\alpha n}\sum_{k=0}^{n}\binom nk\lambda^{k}
= \Bigl(\lambda^{-\alpha}(1 + \lambda)\Bigr)^{n},
\]
dan dengan $\lambda = \frac\alpha{1-\alpha}$:
\[
\lambda^{-\alpha}(1+\lambda)
= \alpha^{-\alpha}(1-\alpha)^{\alpha}\cdot\frac1{1-\alpha}
= \alpha^{-\alpha}(1-\alpha)^{-(1-\alpha)}
= \eu^{H(\alpha)} .
\]
Keoptimalannya: dengan meminimumkan $f(\lambda) = -\alpha\ln\lambda +
\ln(1+\lambda)$ atas $\lambda > 0$, persamaan $f'(\lambda)
= -\frac\alpha\lambda + \frac1{1+\lambda} = 0$ mempunyai satu-satunya
solusi $\lambda = \frac\alpha{1-\alpha}$, yang minimum sebab
$f'' > 0$ --- yakni pilihan pemiringan eksponensial (Chernoff).
Pendamaiannya: menurut pertanyaan 21 satu suku $k =
\lfloor\alpha n\rfloor$ saja sudah berorde
$\eu^{nH(\alpha)}/\sqrt{2\pi\alpha(1-\alpha)n}$, sehingga
\[
\frac{\eu^{nH(\alpha)}}{C\sqrt n} \leq
\sum_{k\leq\alpha n}\binom nk \leq \eu^{nH(\alpha)} :
\]
sehingga laju $H(\alpha)$-nya persis, dan seluruh jumlahnya berbiaya paling banyak
satu faktor $\sqrt n$ atas suku terbesarnya. Lalu dibagi $2^n$,
inilah batas ekor koin adilnya $\P(S_n \leq \alpha n) \leq
\eu^{-n(\ln2 - H(\alpha))}$ --- yakni pemusatan \omterm{def:b3:measure:measure}{ukuran} dalam
satu baris.

\end{solution}
